\documentclass[11pt]{amsart}
\usepackage[T1]{fontenc}
\usepackage{amsmath,amssymb,amsthm,mathtools}
\usepackage[margin=28mm]{geometry}
\usepackage[colorlinks=true,urlcolor=blue]{hyperref}
\newtheorem{theorem}{Theorem}
\newtheorem{lemma}{Lemma}
\newtheorem{corollary}{Corollary}
\theoremstyle{remark}
\newtheorem{remark}{Remark}
\newcommand{\N}{\mathbb N}
\newcommand{\Orb}{\operatorname{Orb}}
\newcommand{\Syr}{\operatorname{Syr}}
\newcommand{\T}{\mathcal T}
\title{A window-excess criterion for polynomial Syracuse orbit heights}
\date{Research note --- 18 November 2026 (conditional; not a solution of Collatz)}
\begin{document}
\maketitle
\begin{abstract}
We deduce a conditional polynomial bound on every positive Syracuse orbit
from two precise inputs: a power-saving bound on the nonnegative window excess
in the first-passage potential decomposition, and a polynomially uniform
lower bound on predecessor densities. The argument uses a finite escape ladder
and a short nonperiodic predecessor selection. It does not prove either
hypothesis, exclude nontrivial cycles, or resolve the Collatz conjecture.
\end{abstract}

\section{Setup}
Let $T:\N\to\N$ be the shortcut $3x+1$ map,
\[
T(x)=\begin{cases}(3x+1)/2,&x\text{ odd},\\x/2,&x\text{ even},\end{cases}
\qquad \Syr(x)=\frac{3x+1}{2^{\nu_2(3x+1)}}\quad(x\text{ odd}).
\]
Let $P(a)=\{w\in\N:T^j(w)=a\text{ for some }j\ge0\}$ and assume
that its logarithmic density $h(a)$ exists, as in the source manuscript.
Set $S_X=\sum_{a\le X}a h(a)$. Fix $\beta\in(0,1/2)$ (the source
manuscript uses $\beta=1/200$), and write
\[
Y_X=X^{1-\beta},\qquad W_X=(Y_X,X]\cap\N,\qquad X_1=X^{1-2\beta}.
\]
For odd $m$, let $\Phi_X(m)$ be the sum, with multiplicities, of
members of $W_X$ appearing in the Syracuse blocks from the first block
through the first odd iterate at most $X_1$. Set $\Phi_X(m)=0$
if there is no such passage. Put $w^\flat=w/2^{\nu_2(w)}$ and
\[
C_X=\lim_{N\to\infty}\frac1{\log N}\sum_{1\le w<N}
\frac{\min\{\Phi_X(w^\flat),H_X\}}{w},\qquad
H_X=X(\log X)^{\varepsilon_0/2}.
\]
Here $\varepsilon_0>0$ is the exponent from the original first-passage
reduction. Assume the density and mean limits of the source manuscript.
We propose the sharper \emph{window excess}
\[
E_X:=\sum_{a\in W_X}a h(a)-C_X
=S_X-S_{Y_X}-C_X.\tag{1}
\]
This is nonnegative: for each starting value the sum of the \emph{distinct}
visited points of $W_X$ dominates $\Phi_X$, since no point can repeat
before a finite first passage below $X_1$; in turn
$\min\{\Phi_X,H_X\}\le\Phi_X$. Therefore
\[
\boxed{S_X=S_{Y_X}+C_X+E_X,\qquad E_X\ge0.}\tag{2}
\]
Unlike the full nonnegative excess $D_X=S_X-C_X=T_X+R_X$ in the source
manuscript, $E_X=D_X-S_{Y_X}$ removes the unavoidable mass of targets
below the upper window. In fact, $E_X$ is the clipping tail plus the
\emph{within-window} part of the return term.

\begin{lemma}[Window excess and the original profile]\label{lem:equivalence}
Assume the source manuscript's calibration
$C_X=\gamma X+o(X)$ with $\gamma=2/\log(4/3)$.
Then $E_X=o(X)$ if and only if $S_X=\gamma X+o(X)$.
\end{lemma}
\begin{proof}
The reverse direction follows immediately from (1), since
$S_{Y_X}=\gamma Y_X+o(Y_X)=o(X)$.
Conversely, if $E_X=o(X)$, then for every $\eta>0$ and sufficiently
large $X$, (2) gives
$S_X\le(\gamma+\eta)X+S_{X^{1-\beta}}$.
Iterate along the scales $x_j=X^{(1-\beta)^j}$ until the first scale
below a fixed cutoff $M>1$. There are $O(\log\log X)$ iterations,
and the sum of all scales except $x_0=X$ is at most
$O(X^{1-\beta}\log\log X)=o(X)$; the terminal $S_{x_j}$ is bounded.
Thus $\limsup S_X/X\le\gamma+\eta$.
The inequality $S_X\ge C_X$ gives $\liminf S_X/X\ge\gamma$.
Let $\eta\downarrow0$.
\end{proof}

\section{A quantitative finite escape ladder}
\begin{lemma}[A high excursion forces window excess]\label{lem:ladder}
Let $a$ be odd and not itself periodic under $T$, with $h(a)>0$.
Suppose $a\le X^{1-2\beta}$, $X^{-\beta}\le1/4$, and an iterate of
$a$ is at least $X$. Then
\[
E_X\ge h(a)\bigl(X-2X^{1-\beta}\bigr)
\ge\frac12h(a)X.\tag{3}
\]
\end{lemma}
\begin{proof}
Write $Y=X^{1-\beta}$ and $b_j=X/2^{j+1}$ for
$j=0,\dots,J$, where $J$ is maximal with $b_J\ge Y$.
Since $a\le X^{1-2\beta}<Y\le b_j$ and the forward orbit of $a$
subsequently reaches at least $X$, there is a first orbit point
$p_j>b_j$ occurring after $a$. Its predecessor is at most $b_j$.
The inequality $T(x)\le(3x+1)/2$ implies, as $b_j\ge Y\ge1$,
\[
b_j<p_j\le\frac32b_j+\frac12\le2b_j.
\]
The disjoint intervals $(b_j,2b_j]$ show that the $p_j$ are distinct.
Moreover they all lie in $W_X$, and
\[
\sum_{j=0}^J p_j>\sum_{j=0}^Jb_j
=X(1-2^{-(J+1)})>X-2Y.
\]
For any $w\in P(a)$, the first Syracuse passage below $X_1$ occurs
no later than the visit to $a$ (or has already occurred at time zero).
Thus every point counted by $\Phi_X(w^\flat)$ occurs before, or at, $a$.
No orbit point after $a$ can be among these points: if an orbit point
before or at $a$ repeated later, $a$ itself would lie on a periodic
orbit. Therefore the distinct $p_j$ lie in the full window-point sum
but are absent from $\Phi_X(w^\flat)$.
The summand defining the window remainder is consequently at least
$\mathbf1_{P(a)}(w)\sum_jp_j$. Taking logarithmic means, and using
clipping only to decrease the first-passage contribution, proves (3).
\end{proof}

\begin{theorem}[Power saving plus predecessor density implies polynomial height]
\label{thm:height}
Assume there are $\delta,K>0$ and $X_*\ge2$ such that
\[
0\le E_X\le K X^{1-\delta}\qquad(X\ge X_*).\tag{4}
\]
Assume further that for some $\kappa,c_*>0$,
\[
h(a)\ge c_*a^{-\kappa}\quad
\text{for all odd }a\ge1\text{ with }3\nmid a.\tag{5}
\]
Then there is $C=C(K,\delta,\beta,c_*,\kappa,X_*)$ such that
for every positive integer $n$,
\[
\boxed{\sup_{j\ge0}T^j(n)\le C n^A,\qquad
A=\max\left\{\frac1{1-2\beta},\frac\kappa\delta\right\}.}\tag{6}
\]
In particular, all positive $T$-orbits are bounded.
\end{theorem}
\begin{proof}
First take an odd nonperiodic $a$ with $3\nmid a$.
If its orbit reaches
$X\ge\max\{X_*,4^{1/\beta},a^{1/(1-2\beta)}\}$,
Lemma~\ref{lem:ladder} and (4) give
$\frac12h(a)X\le K X^{1-\delta}$, hence
\[
X\le\left(\frac{2K}{h(a)}\right)^{1/\delta}
\le\left(\frac{2K}{c_*}\right)^{1/\delta}a^{\kappa/\delta}.
\]
Thus the height of this orbit obeys (6) with $a$ in place of $n$.

Next let $u$ be any odd integer with $3\nmid u$. Choose
$e_0\in\{1,2\}$ with $2^{e_0}u\equiv1\pmod3$ and form three odd
Syracuse predecessors
\[
b_i=\frac{2^{e_0+2i}u-1}{3},\qquad i=0,1,2.
\]
They are positive, each satisfies $\Syr(b_i)=u$, and $b_i<22u$.
Since $b_{i+1}-b_i=2^{e_0+2i}u\not\equiv0\pmod3$ and the differences
have the same nonzero residue modulo $3$, at least two of the $b_i$
are not divisible by $3$. At most one $b_i$ can be periodic: in a
deterministic cycle containing $u$, there is only one predecessor of
$u$ in that cycle. Choose $b$ which is both nonperiodic and
coprime to $3$. The entire forward orbit of $u$ is contained in that
of $b$, and its height is therefore at most $C(22u)^A$.

Finally, for arbitrary $n$, divide by $2$ until reaching an odd $v\le n$.
If $3\nmid v$, put $u=v$; otherwise put $u=\Syr(v)$, which is odd,
not divisible by $3$, and at most $(3v+1)/2\le2n$. The shortcut
orbit from $n$ up to $u$ never exceeds $2n$, and all later terms
belong to the orbit of $u$. Apply the preceding bound to $u\le2n$.
Enlarging $C$ proves (6).
\end{proof}

\begin{corollary}[Conditional polynomial total stopping time]
Under the hypotheses of Theorem~\ref{thm:height}, if the only
positive periodic orbit of $T$ is $1\leftrightarrow2$, then every
positive integer $n$ reaches $1$ within $C' n^A$ shortcut steps.
\end{corollary}
\begin{proof}
There are at most $\lfloor Cn^A\rfloor$ possible positive states in
the orbit. Its first repeated state therefore occurs after at most
$Cn^A+1$ transitions, and belongs to the asserted unique cycle.
\end{proof}

\begin{remark}[Logical status and exponent]
For $\beta=1/200$, $1/(1-2\beta)=100/99$. The source manuscript
establishes $C_X=(2/\log(4/3))X+o(X)$ under its first-passage inputs,
but (4) is a \emph{new, substantially stronger open hypothesis};
indeed even $E_X=o(X)$ is open. Theorem~\ref{thm:height} is not an
unconditional polynomial bound. Recent 2026 non-refereed predecessor
estimates have proposed polynomially uniform bounds (5), including
$\kappa=2+\epsilon$ and a stronger $\kappa=1$ based on additional
joint Syracuse mixing. Their proofs and inputs require independent
validation. With $\kappa=2+\epsilon$ the stated exponent is
$A=\max\{100/99,(2+\epsilon)/\delta\}$; with the stronger estimate
it would be $A=\max\{100/99,1/\delta\}$.
If only $E_X\ll X/(\log X)^p$ can be shown, the same argument gives
at best a stretched-exponential height in $n$, not a polynomial one.
Even (4) by itself does not exclude nontrivial cycles.
\end{remark}

\section{Concrete proof obligation and numerical diagnostic}
Define the finite-orbit observable
\[
A_X(w)=\sum_{a\in\Orb(w)\cap W_X}a,
\qquad
G_X(w)=A_X(w)-\min\{\Phi_X(w^\flat),H_X\}\ge0.
\]
Its logarithmic mean is exactly $E_X$:
\[
E_X=\lim_{N\to\infty}\frac1{\log N}
   \sum_{1\le w<N}\frac{G_X(w)}w.
\]
Therefore the missing analytic input is the \emph{one-sided,
positive-observable} estimate
\[
\boxed{\mathbb E_{\log}G_X(w)\ll X^{1-\delta}}.
\]
It does not involve subtraction of two large measured averages.
A possible route is to write $G_X$ as the clipping tail of the first
passage plus occupation \emph{within} $W_X$ after that passage.
The source manuscript already supplies conditional entrance laws,
two-level return certificates, and a finite-modulus obstruction.
The missing step is a bound uniform in the window scale for the
actual weighted entrance laws, not just a geometric tail in an
independent random-walk surrogate. To claim an unconditional
polynomial height bound, (4) and an independently validated form of
(5) must be proved for the deterministic map.

\end{document}
